\begin{thebibliography}{38}
\bibitem[1]{ref1} Manindra Agrawal, Rohit Gurjar, Arpita Korwar, and Nitin Saxena. Hitting-sets for ROABP and sum of set-multilinear circuits. SIAM J. Comput., 44(3):669–697, 2015. doi:10.1137/140975103.
\bibitem[2]{ref2} Manindra Agrawal and V. Vinay. Arithmetic circuits: A chasm at depth four. In 49th Annual IEEE Symposium on Foundations of Computer Science, FOCS 2008, October 25-28, 2008, Philadelphia, PA, USA, pages 67–75. IEEE Computer Society, 2008. doi:10.1109/FOCS.2008.32. P1
\bibitem[3]{ref3} Miklós Ajtai. $\Sigma^1\_1$-formulae on finite structures. Ann. Pure Appl. Log., 24(1):1–48, 1983. doi:10.1016/0168-0072(83)90038-6.
\bibitem[4]{ref4} Vikraman Arvind and S. Raja. Some lower bound results for set-multilinear arithmetic computations. Chic. J. Theor. Comput. Sci., 2016, 2016. URL: http://cjtcs.cs.uchicago.edu/articles/2016/6/contents.html.
\bibitem[5]{ref5} Peter Bürgisser. Cook’s versus valiant’s hypothesis. Theor. Comput. Sci., 235(1):71–88, 2000. doi:10.1016/S0304-3975(99)00183-8.
\bibitem[6]{ref6} Suryajith Chillara, Christian Engels, Nutan Limaye, and Srikanth Srinivasan. A near-optimal depth-hierarchy theorem for small-depth multilinear circuits. In Mikkel Thorup, editor, 59th IEEE Annual Symposium on Foundations of Computer Science, FOCS 2018, Paris, France, October 7-9, 2018, pages 934–945. IEEE Computer Society, 2018. doi:10.1109/FOCS.2018.00092.
\bibitem[7]{ref7} Suryajith Chillara, Nutan Limaye, and Srikanth Srinivasan. Small-depth multilinear formula lower bounds for iterated matrix multiplication with applications. SIAM J. Comput., 48(1):70– 92, 2019. doi:10.1137/18M1191567.
\bibitem[8]{ref8} Zeev Dvir, Guillaume Malod, Sylvain Perifel, and Amir Yehudayoff. Separating multilinear branching programs and formulas. In Howard J. Karloff and Toniann Pitassi, editors, Proceed- ings of the 44th Symposium on Theory of Computing Conference, STOC 2012, New York, NY, USA, May 19–22, 2012, pages 615–624. ACM, 2012. doi:10.1145/2213977.2214034.
\bibitem[9]{ref9} Michael A. Forbes and Amir Shpilka. Quasipolynomial-time identity testing of non- commutative and read-once oblivious algebraic branching programs. In 54th Annual IEEE Symposium on Foundations of Computer Science, FOCS 2013, 26-29 October, 2013, Berkeley, CA, USA, pages 243–252. IEEE Computer Society, 2013. doi:10.1109/FOCS.2013.34.
\bibitem[10]{ref10} Hervé Fournier, Nutan Limaye, Guillaume Malod, and Srikanth Srinivasan. Lower bounds for depth-4 formulas computing iterated matrix multiplication. SIAM J. Comput., 44(5):1173–1201, 2015. doi:10.1137/140990280.
\bibitem[11]{ref11} Merrick L. Furst, James B. Saxe, and Michael Sipser. Parity, circuits, and the polynomial-time hierarchy. Math. Syst. Theory, 17(1):13–27, 1984. doi:10.1007/BF01744431.
\bibitem[12]{ref12} Zeyu Guo and Rohit Gurjar. Improved explicit hitting-sets for roabps. In Jaroslaw Byrka and Raghu Meka, editors, Approximation, Randomization, and Combinatorial Optimization. Al- gorithms and Techniques, APPROX/RANDOM 2020, August 17-19, 2020, Virtual Conference, volume 176 of LIPIcs, pages 4:1–4:16. Schloss Dagstuhl – Leibniz-Zentrum für Informatik, 2020. doi:10.4230/LIPIcs.APPROX/RANDOM.2020.4.
\bibitem[13]{ref13} Johan Håstad. Almost optimal lower bounds for small depth circuits. In Juris Hartmanis, editor, Proceedings of the 18th Annual ACM Symposium on Theory of Computing, May 28-30, 1986, Berkeley, California, USA, pages 6–20. ACM, 1986. doi:10.1145/12130.12132.
\bibitem[14]{ref14} Pavel Hrubes and Amir Yehudayoff. Homogeneous formulas and symmetric polynomials. Comput. Complex., 20(3):559–578, 2011. doi:10.1007/s00037-011-0007-3.
\bibitem[15]{ref15} Neeraj Kayal, Nutan Limaye, Chandan Saha, and Srikanth Srinivasan. An exponential lower bound for homogeneous depth four arithmetic formulas. SIAM J. Comput., 46(1):307–335, 2017. doi:10.1137/151002423.
\bibitem[16]{ref16} Neeraj Kayal, Vineet Nair, and Chandan Saha. Separation between read-once oblivious algebraic branching programs (roabps) and multilinear depth-three circuits. ACM Trans. Comput. Theory, 12(1):2:1–2:27, 2020. doi:10.1145/3369928.
\bibitem[17]{ref17} Neeraj Kayal, Chandan Saha, and Ramprasad Saptharishi. A super-polynomial lower bound for regular arithmetic formulas. In David B. Shmoys, editor, Symposium on Theory of Computing, STOC 2014, New York, NY, USA, May 31 – June 03, 2014, pages 146–153. ACM, 2014. doi:10.1145/2591796.2591847.
\bibitem[18]{ref18} Neeraj Kayal, Chandan Saha, and Sébastien Tavenas. An almost cubic lower bound for depth three arithmetic circuits. In Ioannis Chatzigiannakis, Michael Mitzenmacher, Yuval Rabani, and Davide Sangiorgi, editors, 43rd International Colloquium on Automata, Languages, and Programming, ICALP 2016, July 11-15, 2016, Rome, Italy, volume 55 of LIPIcs, pages 33:1–33:15. Schloss Dagstuhl – Leibniz-Zentrum für Informatik, 2016. doi:10.4230/LIPIcs.ICALP.2016.33.
\bibitem[19]{ref19} Neeraj Kayal, Chandan Saha, and Sébastien Tavenas. On the size of homogeneous and of depth-four formulas with low individual degree. Theory Comput., 14(1):1–46, 2018. doi:10.4086/toc.2018.v014a016.
\bibitem[20]{ref20} Pascal Koiran. Arithmetic circuits: The chasm at depth four gets wider. Theor. Comput. Sci., 448:56–65, 2012. doi:10.1016/j.tcs.2012.03.041.
\bibitem[21]{ref21} Mrinal Kumar and Shubhangi Saraf. On the power of homogeneous depth 4 arithmetic circuits. SIAM J. Comput., 46(1):336–387, 2017. doi:10.1137/140999335.
\bibitem[22]{ref22} Deepanshu Kush and Shubhangi Saraf. Improved low-depth set-multilinear circuit lower bounds. In Shachar Lovett, editor, 37th Computational Complexity Conference, CCC 2022, July 20-23, 2022, Philadelphia, PA, USA, volume 234 of LIPIcs, pages 38:1–38:16. Schloss Dagstuhl – Leibniz-Zentrum für Informatik, 2022. doi:10.4230/LIPIcs.CCC.2022.38.
\bibitem[23]{ref23} Nutan Limaye, Srikanth Srinivasan, and Sébastien Tavenas. Superpolynomial lower bounds against low-depth algebraic circuits. In 62nd IEEE Annual Symposium on Foundations of Computer Science, FOCS 2021, Denver, CO, USA, February 7-10, 2022, pages 804–814. IEEE, 2021. doi:10.1109/FOCS52979.2021.00083.
\bibitem[24]{ref24} Noam Nisan and Avi Wigderson. Hardness vs randomness. J. Comput. Syst. Sci., 49(2):149–167, 1994. doi:10.1016/S0022-0000(05)80043-1.
\bibitem[25]{ref25} Noam Nisan and Avi Wigderson. Lower bounds on arithmetic circuits via partial derivatives. Comput. Complex., 6(3):217–234, 1997. doi:10.1007/BF01294256.
\bibitem[26]{ref26} Ran Raz. Separation of multilinear circuit and formula size. Theory Comput., 2(6):121–135, 2006. doi:10.4086/toc.2006.v002a006.
\bibitem[27]{ref27} Ran Raz. Multi-linear formulas for permanent and determinant are of super-polynomial size. J. ACM, 56(2):8:1–8:17, 2009. doi:10.1145/1502793.1502797.
\bibitem[28]{ref28} Ran Raz. Tensor-rank and lower bounds for arithmetic formulas. J. ACM, 60(6):40:1–40:15, 2013. doi:10.1145/2535928.
\bibitem[29]{ref29} Ran Raz and Amir Yehudayoff. Balancing syntactically multilinear arithmetic circuits. Comput. Complex., 17(4):515–535, 2008. doi:10.1007/s00037-008-0254-0.
\bibitem[30]{ref30} Ran Raz and Amir Yehudayoff. Lower bounds and separations for constant depth multilinear circuits. Comput. Complex., 18(2):171–207, 2009. doi:10.1007/s00037-009-0270-8.
\bibitem[31]{ref31} Alexander A. Razborov. Lower bounds on the size of bounded depth circuits over a complete basis with logical addition. Mathematical notes of the Academy of Sciences of the USSR, 41:333–338, 1987.
\bibitem[32]{ref32} Ramprasad Saptharishi. A survey of lower bounds in arithmetic circuit complexity. Github Survey, 2015. URL: https://github.com/dasarpmar/lowerbounds-survey.
\bibitem[33]{ref33} Amir Shpilka and Amir Yehudayoff. Arithmetic circuits: A survey of recent results and open questions. Found. Trends Theor. Comput. Sci., 5(3-4):207–388, 2010. doi:10.1561/0400000039.
\bibitem[34]{ref34} Roman Smolensky. Algebraic methods in the theory of lower bounds for boolean circuit complexity. In Alfred V. Aho, editor, Proceedings of the 19th Annual ACM Symposium on Theory of Computing, 1987, New York, New York, USA, pages 77–82. ACM, 1987. doi:10.1145/28395.28404.
\bibitem[35]{ref35} Sébastien Tavenas. Improved bounds for reduction to depth 4 and depth 3. Inf. Comput., 240:2–11, 2015. doi:10.1016/j.ic.2014.09.004.
\bibitem[36]{ref36} Sébastien Tavenas, Nutan Limaye, and Srikanth Srinivasan. Set-multilinear and non- commutative formula lower bounds for iterated matrix multiplication. In Stefano Le- onardi and Anupam Gupta, editors, STOC ’22: 54th Annual ACM SIGACT Symposium on Theory of Computing, Rome, Italy, June 20–24, 2022, pages 416–425. ACM, 2022. doi:10.1145/3519935.3520044.
\bibitem[37]{ref37} Leslie G. Valiant, Sven Skyum, Stuart J. Berkowitz, and Charles Rackoff. Fast parallel computation of polynomials using few processors. SIAM J. Comput., 12:641–644, 1983.
\bibitem[38]{ref38} Andrew Chi-Chih Yao. Separating the polynomial-time hierarchy by oracles (preliminary version). In 26th Annual Symposium on Foundations of Computer Science, Portland, Oregon, USA, 21-23 October 1985, pages 1–10. IEEE Computer Society, 1985. doi:10.1109/SFCS.1985.49.
\end{thebibliography}
